\documentclass[12pt,a4paper]{article}
\usepackage[UTF8]{ctex}
\usepackage{amsmath,amssymb}
\usepackage{geometry}
\usepackage{booktabs}
\usepackage{array}
\usepackage{xcolor}
\usepackage{tcolorbox}

\geometry{left=2cm,right=2cm,top=2cm,bottom=2cm}

\title{数学错题分析报告}
\author{WrongMath OCR 识别}
\date{\today}

\begin{document}

\maketitle

\section{计算题部分}

\subsection{直接写出得数}

\begin{table}[h]
\centering
\begin{tabular}{|c|c|c|c|}
\hline
\textbf{题号} & \textbf{题目} & \textbf{答案} & \textbf{批改} \\
\hline
1 & $12 \times 80$ & 960 & \textcolor{green}{\checkmark} \\
\hline
2 & $96 \div 24$ & 4 & \textcolor{green}{\checkmark} \\
\hline
3 & $840 - 240 \div 6$ & 800 & \textcolor{green}{\checkmark} \\
\hline
4 & $500 \div 70$ & $7 \cdots 10$ & \textcolor{green}{\checkmark} \\
\hline
5 & $400 \div 50$ & 8 & \textcolor{green}{\checkmark} \\
\hline
6 & $12 \times 101$ & 1212 & \textcolor{green}{\checkmark} \\
\hline
7 & $1250 \div (25 \times 5)$ & 10 & \textcolor{green}{\checkmark} \\
\hline
8 & $(100) + 100 = (20) \times 10$ & 100, 20 & \textcolor{green}{\checkmark} \\
\hline
\end{tabular}
\end{table}

\subsection{探究计算方法}

\textbf{应用题：}每套书有 14 本，李老师买了 12 套，一共买了多少本？

\textbf{方法一（分拆为 $4 \times 3$）：}
\begin{align*}
14 \times 12 &= 14 \times 3 \times 4 \\
&= 42 \times 4 \\
&= 168
\end{align*}

\textbf{方法二（分拆为 $10 + 2$）：}
\begin{align*}
14 \times 12 &= 14 \times 10 + 14 \times 2 \\
&= 140 + 28 \\
&= 168
\end{align*}

\subsection{列竖式计算}

\textbf{(1)} $320 \times 45 = 14400$ \quad \textcolor{green}{\checkmark}

\textbf{(2)} $4200 \div 28 = 150$ \quad \textcolor{green}{\checkmark}

\textbf{(3)} $2\square06 \div 24$ \quad \textcolor{red}{\ding{55} \textbf{错题，扣14分}}

\begin{tcolorbox}[colback=red!5!white,colframe=red!75!black,title=错题详情]
\textbf{题目要求：}先在$\square$里填数，使商中间有0，再列竖式计算。

\textbf{学生答案：}填了4，即 $2406 \div 24$

\textbf{学生竖式：}
\[
\begin{array}{r}
      100 \\
\cline{2-2}
24 \overline{)2406} \\
     24 \\
\cline{2-2}
       06
\end{array}
\]

\textbf{错误分析：}
\begin{itemize}
    \item 学生计算 $2406 \div 24 = 100$ 余 6
    \item 但题目要求"使商中间有0"
    \item 正确答案应填 \textbf{4} 或 \textbf{8}，使商为三位数且中间为0
    \item 正确计算：$2406 \div 24 = 100 \cdots 6$ 或 $2806 \div 24$
\end{itemize}
\end{tcolorbox}

\subsection{递等式计算}

\begin{align*}
600 \div 24 + 76 &= 25 + 76 = 101 \quad \textcolor{green}{\checkmark} \\
2105 - 105 \times 16 &= 2105 - 1680 = 425 \quad \textcolor{green}{\checkmark} \\
(28 + 55) \times 22 &= 83 \times 22 = 1826 \quad \textcolor{green}{\checkmark} \\
3380 \div (650 \div 50) &= 3380 \div 13 = 260 \quad \textcolor{green}{\checkmark}
\end{align*}

\section{错题汇总}

\begin{table}[h]
\centering
\begin{tabular}{|c|c|c|c|}
\hline
\textbf{题号} & \textbf{题目} & \textbf{错误类型} & \textbf{扣分} \\
\hline
3-(3) & $2\square06 \div 24$ & 填空及竖式计算错误 & -14分 \\
\hline
\end{tabular}
\end{table}

\section{知识点分析}

\subsection{错题涉及知识点}
\begin{itemize}
    \item \textbf{除法竖式计算}：三位数除以两位数
    \item \textbf{商的性质}：理解"商中间有0"的含义
    \item \textbf{验算方法}：乘法验算除法
\end{itemize}

\subsection{建议练习方向}
\begin{enumerate}
    \item 加强对"商中间有0"概念的理解
    \item 练习构造满足特定条件的除法算式
    \item 巩固除法验算方法
\end{enumerate}

\end{document}